\section{Algorithm: Multi-scale MCMC Option B}
\label{sec:algorithm}

\subsection{Setup and Notation}

Let $G = (V, E)$ be the census-tract adjacency graph for a state, with
$|V| = n$ tracts.
Each tract $v$ has population $\mathrm{pop}(v)$ and a GEOID string
$\mathrm{geoid}(v)$ of 11 digits (2 state + 3 county + 6 tract).
The county FIPS code of tract $v$ is
\[
  \mathrm{county}(v) = \mathrm{geoid}(v)[0{:}5],
\]
the first five characters of the GEOID.
Let $C = \{\mathrm{county}(v) : v \in V\}$ be the set of distinct county
FIPS codes, $|C| = n_C$.

The ideal fine population is $P_{\mathrm{ideal}} = \sum_v \mathrm{pop}(v) / k$;
the fine tolerance is $\epsilon_f$ (default 0.005).
The CLI wrapper uses coarse tolerance $\epsilon_c = 3 \epsilon_f$ to reflect
the lower precision of county-level population accounting and to avoid rejecting
nearly every coarse proposal before the fine-level rebalance step has a chance
to repair it.
The library configuration substrate still records a fixed-alpha default of
$2\epsilon_f$, while the adaptive configuration and current CLI wrapper use
$3\epsilon_f$.
This mismatch is an implementation-boundary item: reports should record the
actual tolerance multiplier used by the run rather than relying on a paper
default.

\subsection{County Adjacency Derivation}

\begin{proposition}[County adjacency from the run's tract graph]
\label{prop:county-adj-exact}
Define the county adjacency graph $G_C = (C, E_C)$ by
\[
  (c_1, c_2) \in E_C
  \iff c_1 \neq c_2
  \text{ and }
  \exists\, (v_1, v_2) \in E \text{ with }
  \mathrm{county}(v_1) = c_1,\;
  \mathrm{county}(v_2) = c_2.
\]
This derivation is exact relative to the tract adjacency graph used by the run:
two coarse county nodes are adjacent exactly when the run's fine graph contains
at least one cross-county tract edge.
The derivation requires one pass over the tract edge list and runs in $O(|E|)$
time.
\end{proposition}

\begin{proof}
The algorithm scans the fine adjacency list and inserts a county edge whenever
the endpoints' county prefixes differ. Therefore every inserted coarse edge has
a witnessing fine edge, and every witnessing fine edge is represented in the
coarse graph. Geographic adjacency not represented in the input tract graph
cannot be recovered by this derivation.
\end{proof}

\begin{remark}[Water-boundary caveat]
Proposition~3.1 assumes that two geographically adjacent counties share at least one cross-county tract edge in the TIGER/Line adjacency graph. This holds for land-boundary adjacency. However, some counties are geographically adjacent only across water bodies (e.g., two counties on opposite banks of a river whose tract boundaries do not cross). In TIGER/Line adjacency graphs built without water-crossing bridges, such county pairs may be absent from the derived county adjacency graph---producing a county graph that is slightly sparser than geographic adjacency.

For redistricting purposes, this is typically acceptable: federal congressional districts are not required to cross water bodies, and the missing county adjacency represents a pair that would be difficult to connect in practice. Users in states with significant water-separated county pairs (e.g., Texas Gulf Coast, Louisiana bayou regions) should verify the derived county adjacency against geographic county maps.
\end{remark}

The derivation scans each edge $(v_1, v_2) \in E$; if
$\mathrm{county}(v_1) \neq \mathrm{county}(v_2)$, the pair is added to
$E_C$.
For Texas, $|E| \approx 15{,}600$ tract edges yield $|E_C| \approx 220$
county edges from $|C| = 254$ counties.

\subsection{County Population and Initial Coarse Plan}

The county population is
$\mathrm{pop}(c) = \sum_{v \in f^{-1}(c)} \mathrm{pop}(v)$.
The model-level coarse plan $\pi_C^{(0)} : C \to [k]$ can be derived from the
initial fine plan $\pi^{(0)} : V \to [k]$ by \emph{plurality projection}:
\[
  \pi_C^{(0)}(c) = \arg\max_{d \in [k]}\;
    \mathrm{pop}\!\left(\{v \in f^{-1}(c) : \pi^{(0)}(v) = d\}\right).
\]
That is, each county is assigned to the district that contains the plurality
of its population. The current \texttt{run\_multiscale} wrapper instead
initializes and syncs the coarse assignment by iterating
\texttt{fine\_to\_coarse} and assigning the observed fine district label to the
coarse slot. This is a simpler implementation rule than the model-level
plurality projection above and should be made explicit in audit artifacts.

\subsection{Rebalance Subroutine}

After a coarse move proposes a new $\pi_C'$, the fine plan is updated by
assigning each tract to the district of its county:
$\pi'(v) = \pi_C'(\mathrm{county}(v))$.
This projection may violate fine population balance because counties are not
perfectly equal in population.
The \emph{rebalance subroutine} restores balance by performing boundary swaps:

\begin{enumerate}
  \item Identify all districts $d$ with $|\mathrm{pop}(d) - P_{\mathrm{ideal}}|
        > \epsilon_f \cdot P_{\mathrm{ideal}}$.
  \item For each over-populated district, find boundary tracts adjacent to
        under-populated districts and move the tract with the best
        population-balance improvement.
  \item Repeat up to 200 iterations.
  \item If any district remains unbalanced after 200 iterations, \emph{reject}
        the coarse move entirely (revert to the previous fine plan $\pi$).
\end{enumerate}

The 200-iteration limit is a hard contract: it prevents the rebalance from
consuming unbounded time while still accommodating the typical imbalances
introduced by county-level projection.
In practice, for Texas, rebalance completes within 5--20 iterations on
accepted coarse moves; failure occurs when a coarse move produces a coarse
district that contains an entire county whose population far exceeds
$P_{\mathrm{ideal}}$.

In the atlas framing, a multiscale move is a candidate coarse/fine transition
with a visible gate.
Coarse moves propose on the county graph, project back to tracts, and then
either pass the rebalance gate or are rejected with the previous fine plan
restored.
Fine moves stay on the tract graph.
The transcript should preserve which scale proposed the move, whether
projection and rebalance succeeded, and what status can be claimed for the
accepted plan.

\subsection{Complete Algorithm}

\paragraph{Algorithm 1. \ms{} Option B as implemented by \texttt{run\_multiscale}.}
\begin{algorithmic}[1]
\Require Tract graph $G$, populations $\{\mathrm{pop}(v)\}$, GEOIDs,
         initial plan $\pi^{(0)}$, districts $k$,
         fine tolerance $\epsilon_f$, coarse factor 3,
         coarsening probability $\alpha$, total steps $T$,
         percentile $p \in [0,1]$, base seed $b$
\Ensure  Selected plan $\pi^*$

\State $\epsilon_c \gets 3 \epsilon_f$
\State $f(v) \gets \mathrm{geoid}(v)[0{:}5]$ for all $v \in V$
        \Comment{tract $\to$ county map}
\State $G_C, \{\mathrm{pop}(c)\} \gets \textsc{DeriveCountyGraph}(G, f)$
        \Comment{$O(|E|)$, Prop.~\ref{prop:county-adj-exact}}
\State $\pi_C \gets \textsc{SyncCoarseFromFine}(\pi^{(0)}, f)$
\State $\pi \gets \pi^{(0)}$
\State $\mathrm{visited} \gets [(\ec(\pi),\,0,\, \pi)]$

\For{$s \gets 1$ \textbf{to} $T$}
  \State $u \gets \mathrm{SHA256}(\texttt{"MSC\_STEP\_"} \,\|\, s^{\mathrm{u64le}}
         \,\|\, \texttt{"\_"} \,\|\, 0^{\mathrm{u32le}} \,\|\, \texttt{"\_"} \,\|\, b^{\mathrm{u64le}})[0{:}8]$
         as $\mathrm{SmallRng}$
  \If{$u.\mathrm{gen\_range}(0.0, 1.0) < \alpha$}
    \Comment{Coarse move}
    \State $\pi_C' \gets \recom{}.\mathrm{step}(G_C, \{\mathrm{pop}(c)\}, \pi_C, \epsilon_c, u)$
    \State $\pi' \gets \textsc{Project}(\pi_C', f)$
            \Comment{assign each tract to its county's district}
    \State $\mathrm{ok} \gets \textsc{Rebalance}(\pi', \epsilon_f, \text{max}=200)$
    \If{$\mathrm{ok}$}
      \State $\pi \gets \pi'$;\quad $\pi_C \gets \textsc{SyncCoarseFromFine}(\pi, f)$
    \EndIf
    \Comment{else reject: keep $\pi$, $\pi_C$ unchanged}
  \Else
    \Comment{Fine move}
    \State $\pi \gets \recom{}.\mathrm{step}(G, \{\mathrm{pop}(v)\}, \pi, \epsilon_f, u)$
    \State $\pi_C \gets \textsc{SyncCoarseFromFine}(\pi, f)$
            \Comment{keep coarse plan consistent}
  \EndIf
  \State $\mathrm{visited}.\mathrm{push}((\ec(\pi), s, \pi))$
\EndFor

\State Sort $\mathrm{visited}$ by $(\ec,\mathrm{step})$ ascending
\State $r \gets \min(\lfloor p \cdot |\mathrm{visited}|\rfloor, |\mathrm{visited}|-1)$
\State \Return $\mathrm{visited}[r].\mathrm{plan}$
\end{algorithmic}

\paragraph{Deterministic seed schedule.}
Each step $s$ derives its RNG from a SHA-256 hash of the step index and
the base seed:
\[
  \mathrm{step\_seed}(s, b)
  = \mathrm{SHA256}\!\left(\texttt{"MSC\_STEP\_"}
    \,\|\, s^{\mathrm{u64le}}
    \,\|\, \texttt{"\_"}
    \,\|\, 0^{\mathrm{u32le}}
    \,\|\, \texttt{"\_"}
    \,\|\, b^{\mathrm{u64le}}\right)[0{:}8]
    \;\text{as}\; u64.
\]
This ensures that the sequence of coarse/fine decisions---and the outcomes
of each step---is fully determined by $b$ and $T$, enabling exact
reproduction of any run from the manifest entry alone.

\paragraph{CLI invocation.}
\begin{verbatim}
bisect state --state TX --year 2020 \
    --search multiscale \
    --multiscale-steps 5000 \
    --multiscale-alpha 0.3 \
    --multiscale-fine tract \
    --multiscale-coarse county \
    --percentile 0.0
\end{verbatim}

\subsection{Correctness Properties}

\begin{proposition}[Fine plan validity]
\label{prop:validity}
Every fine plan $\pi$ in the \texttt{visited} list that passes the current
validity gates is population-balanced within $\epsilon_f$; the implementation
intends accepted coarse moves to be rebalanced before they enter the selected
plan pool.
\end{proposition}

\begin{proof}
The initial plan $\pi^{(0)}$ is valid by construction from \cs{}/METIS.
Fine \recom{} steps preserve validity (the standard \recom{} invariant).
Coarse moves are accepted only if \textsc{Rebalance} returns \texttt{ok},
which checks $|\mathrm{pop}(d) - P_{\mathrm{ideal}}| \leq \epsilon_f \cdot
P_{\mathrm{ideal}}$ for all $d$ after rebalancing.
Connectivity of fine districts after a coarse move: each tract is assigned
to the district of its county; within a county, the district assignment is
uniform, so connectivity of the county in $G_C$ is necessary but not
sufficient for connectivity of the fine district.
The rebalance swaps maintain a connected fine district by restricting to
boundary-tract swaps that do not disconnect the donor district (checked via
BFS, as in the \flip{} chain).
\end{proof}

\begin{remark}[Stationarity]
\label{rem:stationarity}
Like standard \recom{}, the current \ms{} wrapper does not have a fully characterised stationary
distribution.
The coarse moves introduce an additional source of non-reversibility: the
projection and rebalance operations are not symmetric in forward and reverse
directions.
For practical redistricting ensemble analysis, this is acceptable---the goal
is faster mixing to explore the plan space, not exact sampling from a specific
distribution.
Applications requiring stronger distributional guarantees need a separate
acceptance accounting for the coarse projection/rebalance move. A ForestReCom
fine chain~\citep{dellaLibera2026forestRecom} would improve only the fine-level
component unless the coarse component is also Metropolized.
\end{remark}
